% =====================================================================
%  Chapter 10 — The State of Play, Part I
% =====================================================================

\part{The State of Play}
\markboth{Part IV}{The State of Play}

\chapter{The Independent Compilations}
\label{ch:compilations}

The mirror's README said, on the morning of 20 November 2026, that no
independent build had been published, and that the code should be treated as
``a claimed proof until it has been compiled and checked''. By lunchtime, the
first half of that sentence had been overtaken by events. This chapter
reports what happened, and what it does and does not establish.

\section{What was reported}

Quanta Magazine's piece, by Konstantin Kakaes, is the fullest account:

\begin{quote}
\small
``Several independent efforts have now successfully compiled Alpöge's proof.
\textbf{Andrew Sutherland at the Massachusetts Institute of Technology was the
first}\idx{Sutherland, Andrew}; he finished shortly before noon. A few hours later, \textbf{Simon DeDeo
at Carnegie Mellon University}\idx{DeDeo, Simon} also succeeded, and not long after that
Sutherland compiled it using a variant of Lean\idx{variant of Lean}.''
\end{quote}

Three things in that sentence deserve comment.

\begin{description}[style=nextline,leftmargin=0pt,font=\sffamily\bfseries\small,
  itemsep=5pt]
\item[\textbf{Two independent compilers.}] Sutherland at MIT and DeDeo at CMU
are two groups, two machines, two toolchain setups. Agreement between them is
evidence that the build is reproducible rather than a local accident.
\item[\textbf{A variant of Lean.}] The final clause --- ``Sutherland compiled
it using a variant of Lean'' --- is the strongest single piece of evidence in
the story, and it is easy to read past. If the development compiles under a
\emph{modified} Lean, then the modification is doing something specific, and
the most likely something is an extension that makes the build tractable.
\item[\textbf{The order of events.}] ``Shortly before noon'' on 20 November.
The mirror's README was written that morning. The status line in it was true
when written and became obsolete within hours.
\end{description}

\section{What this establishes}

A successful independent compilation establishes the following, and it is a
substantial list:

\begin{itemize}
\item The development is \emph{complete} in the sense that matters: it has a
consistent, terminating development in which the headline theorems exist and
are accepted by the Lean kernel.
\item It is \emph{reproducible}: two groups, working independently, got the
same result.
\item It is \emph{pinned}: the toolchain and the Mathlib revision are fixed, so
the build refers to a determinate mathematical library.
\item The axiom question is, in practice, settled in the affirmative: the
compilers necessarily ran the axiom check, or would have had to.
\end{itemize}

It is a genuine milestone, and the book says so.

\section{What this does not establish}

It does not establish that the mathematics is right. This is not a criticism of
the compilers; it is a statement about what compilation is.

\begin{quote}
``Lean, besides being a mathematical language, is also a computer programming
language, and \textbf{can itself contain bugs. It can also be used to play
sneaky tricks by, for example, redefining symbols used in a proof to make
something trivially true.} The first step in verifying a proof in Lean is to
make sure it `compiles.'\,''
\qquad --- Quanta Magazine, 2026-11-20
\end{quote}

The compilers did the first step. The second step --- checking --- is a
different activity, performed by a different kind of person, and it has not
happened.

\section{The build attempt made for this book}

For completeness, and because a book that criticises others for not compiling
ought to compile, an independent build was attempted in the course of writing
this one.

\begin{description}[style=nextline,leftmargin=0pt,font=\sffamily\bfseries\small,
  itemsep=5pt]
\item[The toolchain.] \texttt{elan} was installed and the pinned toolchain
\code{leanprover/lean4:v4.33.1} was fetched.
\item[The library.] Mathlib was checked out and confirmed at exactly
\code{0df444a360eaa60ab8c11dca51a86af692955474}, the revision pinned by the
submission's \texttt{lakefile.toml}. The prebuilt artefacts were retrieved from
the community cache --- 8{,}690 of them.
\item[The build.] \texttt{lake build Main} was launched against the project's
own sources.
\end{description}

The build was in progress --- compiling the submission's own files --- when the
writing of this book reached the point where its findings had to be written
down. The machine available for it has two cores and three gigabytes of RAM;
the mirror's own README warns that this is ``a very large build ($\approx 17$
M lines on top of Mathlib)'' for which one should ``use a machine with many
cores and plenty of RAM, and expect it to run for a long time''.

\begin{seamnote}{Reporting the outcome honestly}
This book's build did not complete within the time available, and that is
reported as what it is: an incomplete verification attempt, not a failed one
and not a successful one. The log is retained.

The finding it does contribute is narrow and worth stating: the submission
builds far enough to be compiling the project's own modules, with Mathlib at
the pinned revision satisfied entirely from cache, with no error. That is
consistent with the two reported independent compilations, and it adds
nothing to them.

It also contributes a practical datum that the mirror's README lacked: the
build is tractable with the community's prebuilt Mathlib cache, and the
binding constraint is the size of the submission's own source, not the
dependency. A reader with a large machine can complete it.
\end{seamnote}

\section{Why ``it compiles'' is the beginning and not the end}

There is a temptation, on reading a story like this one, to treat compilation
as the decisive event: the machine said yes, so the theorem is proved. It is
worth resisting that, and the reason is not scepticism about the compilers. It
is that the community has been here before.

The formalisation community's own history contains cases where a development
compiled cleanly, passed the axiom check, had its statements reviewed --- and
was then found to be wrong. The most instructive recent example is recorded in
OpenAI's own project history: a formalised result was withdrawn in October 2026
with a stated sign error\idx{sign error (withdrawn formalisation)}. The development compiled. The axiom list was clean.
The failure was in the mathematics of the statement, and no compiler can see
it.

\begin{quote}
This is the ordinary condition of the subject, and it is why the phrase
``until it has been compiled \emph{and checked}'' in the mirror's README has
two clauses rather than one. The second clause has not happened.
\end{quote}

\section{What a completed verification would look like}

If someone were to complete the full protocol --- compile, replay the kernel
independently, and perform a full statement review --- what would they have
established? This is worth spelling out, because it defines the target that
Chapter~14 aims at.

They would have established that the Riemann hypothesis is a theorem of Lean
4, in the sense that the development contains a proof of Mathlib's
\texttt{RiemannHypothesis} whose only dependencies are the three standard
axioms and whose statements, on review, match the classical claims.

That would be an enormous result. It would settle the problem. And it would
still require the community to check the review --- one human reading of
175{,}000 files is itself a document with error bars, and the standard practice
in the field, as the tooling of the Lean FRO shows, is to have the review
independently reproduced.

\begin{keytakeaways}

\item Two independent compilations succeeded on 20 November 2026, by
Andrew Sutherland at MIT and Simon DeDeo at CMU, and a third using a variant
of Lean.
\item Compilation establishes completeness, reproducibility, pinning and the
axiom question. It does not establish that the statements mean what the claim
says they mean.
\item The build attempted for this book reached the stage of compiling the
submission's own modules with Mathlib satisfied from cache, on a
two-core machine, and did not complete within the time available.
\item The community's own history contains a case of a cleanly compiling
development withdrawn for a mathematical error. That is why ``compiled and
checked'' has two clauses.
\end{keytakeaways}
